\begin{thebibliography}{99}
\footnotesize\itemsep=0pt

\bibitem{B}
Bass H., A nontriangular action of {${\mathbb G}_{a}$} on {${\mathbb A}^{3}$},
 \href{https://doi.org/10.1016/0022-4049(84)90019-7}{\textit{J.~Pure Appl. Algebra}} \textbf{33} (1984), 1--5.

\bibitem{Co}
Cohn P.M., Free rings and their relations, 2nd ed., \textit{London Mathematical Society
 Monographs}, Vol.~19, Academic Press, Inc., London, 1985.

\bibitem{Cz1}
Czerniakiewicz A.J., Automorphisms of a free associative algebra of
 rank~{$2$}.~{I}, \href{https://doi.org/10.2307/1995814}{\textit{Trans. Amer. Math. Soc.}} \textbf{160} (1971),
 393--401.

\bibitem{Cz2}
Czerniakiewicz A.J., Automorphisms of a free associative algebra of
 rank~{$2$}.~{II}, \href{https://doi.org/10.2307/1996385}{\textit{Trans. Amer. Math. Soc.}} \textbf{171} (1972),
 309--315.

\bibitem{Jo}
de~W.~Jooste T., Primitive derivations in free associative algebras,
 \href{https://doi.org/10.1007/BF01214786}{\textit{Math.~Z.}} \textbf{164} (1978), 15--23.

\bibitem{D1}
Drensky V., Free algebras and {PI}-algebras. {G}raduate course in algebra,
 Springer-Verlag, Singapore, 2000.

\bibitem{DG}
Drensky V., Gupta C.K., Constants of {W}eitzenb\"{o}ck derivations and
 invariants of unipotent transformations acting on relatively free algebras,
 \href{https://doi.org/10.1016/j.jalgebra.2005.07.004}{\textit{J.~Algebra}} \textbf{292} (2005), 393--428, \href{https://arxiv.org/abs/math.RA/0412399}{arXiv:math.RA/0412399}.

\bibitem{Fa}
Falk G., Konstanzelemente in {R}ingen mit {D}ifferentiation, \href{https://doi.org/10.1007/BF01343560}{\textit{Math.
 Ann.}} \textbf{124} (1952), 182--186.

\bibitem{Fr1}
Freudenburg G., A survey of counterexamples to {H}ilbert's fourteenth problem,
 \textit{Serdica Math.~J.} \textbf{27} (2001), 171--192.

\bibitem{Fr2}
Freudenburg G., Algebraic theory of locally nilpotent derivations,
 \textit{Encyclopaedia of Mathematical Sciences}, Vol.~136, \href{https://doi.org/10.1007/978-3-540-29523-5}{Springer-Verlag},
 Berlin, 2006.

\bibitem{J}
Jung H.W.E., \"{U}ber ganze birationale {T}ransformationen der {E}bene,
 \href{https://doi.org/10.1515/crll.1942.184.161}{\textit{J.~Reine Angew. Math.}} \textbf{184} (1942), 161--174.

\bibitem{Kh}
Kharchenko V.K., Algebras of invariants of free algebras, \href{https://doi.org/10.1007/BF01674783}{\textit{Algebra
 Logic}} \textbf{17} (1978), 316--321, translation of \textit{Algebra i Logika} \textbf{17} (1978), 478--487.

\bibitem{L}
Lane D.R., Free algebras of rank two and their automorphisms, Ph.D.~Thesis,
 {B}edford College, London, 1976.

\bibitem{ML}
Makar-Limanov L., Automorphisms of a free algebra with two generators,
 \href{https://doi.org/10.1007/BF01075252}{\textit{Funct. Anal. Appl.}} \textbf{4} (1970), 262--264,
 translation of \textit{Funk. Analiz i ego Prilozh.} \textbf{4} (1970),  no.~3, 107--108.

\bibitem{ML2}
Makar-Limanov L., Locally nilpotent derivations, a new ring invariant and applications,
{L}ecture notes, available at \url{http://www.math.wayne.edu/~lml/lmlnotes.pdf}.

\bibitem{MLTU}
Makar-Limanov L., Turusbekova U., Umirbaev U., Automorphisms and derivations of
 free {P}oisson algebras in two variables, \href{https://doi.org/10.1016/j.jalgebra.2008.01.005}{\textit{J.~Algebra}} \textbf{322}
 (2009), 3318--3330, \href{https://arxiv.org/abs/0708.1148}{arXiv:0708.1148}.

\bibitem{Na}
Nagata M., On automorphism group of {$k[x,y]$}, \textit{Lectures in Math.},
 Vol.~5, Kyoto University, Kinokuniya Book-Store Co., Ltd., Tokyo, 1972.

\bibitem{No1}
Nowicki A., Polynomial derivations and their rings of constants, Uniwersytet
 Miko{\l}aja Kopernika, Toru\'{n}, 1994.

\bibitem{No2}
Nowicki A., The fourteenth problem of {H}ilbert for polynomial derivations, in
 Differential {G}alois {T}heory ({B}\c{e}dlewo, 2001), \textit{Banach Center
 Publ.}, Vol.~58, \href{https://doi.org/10.4064/bc58-0-13}{Polish Acad. Sci. Inst. Math.}, Warsaw, 2002, 177--188.

\bibitem{R}
Rentschler R., Op\'{e}rations du groupe additif sur le plan affine,
 \textit{C.~R.~Acad. Sci. Paris S\'{e}r.~A-B} \textbf{267} (1968), 384--387.

\bibitem{Se}
Serre J.-P., Trees, \textit{Springer Monographs in Mathematics}, Springer-Verlag, Berlin,
 2003, translation of Arbres, amalgames, ${\rm SL}_2$,
\textit{Ast\'erisque}, Vol.~46,
Soci\'et\'e Math\'ematique de France, Paris, 1977.

\bibitem{Ses}
Seshadri C.S., On a theorem of {W}eitzenb\"{o}ck in invariant theory,
 \href{https://doi.org/10.1215/kjm/1250525012}{\textit{J.~Math. Kyoto Univ.}} \textbf{1} (1962), 403--409.

\bibitem{SU1}
Shestakov I.P., Umirbaev U.U., Poisson brackets and two-generated subalgebras
 of rings of polynomials, \href{https://doi.org/10.1090/S0894-0347-03-00438-7}{\textit{J.~Amer. Math. Soc.}} \textbf{17} (2004),
 181--196.

\bibitem{SU2}
Shestakov I.P., Umirbaev U.U., The tame and the wild automorphisms of
 polynomial rings in three variables, \href{https://doi.org/10.1090/S0894-0347-03-00440-5}{\textit{J.~Amer. Math. Soc.}} \textbf{17}
 (2004), 197--227.

\bibitem{Sm}
Smith M.K., Stably tame automorphisms, \href{https://doi.org/10.1016/0022-4049(89)90158-8}{\textit{J.~Pure Appl. Algebra}}
 \textbf{58} (1989), 209--212.

\bibitem{Sp}
Specht W., Gesetze in {R}ingen.~{I}, \href{https://doi.org/10.1007/BF02230710}{\textit{Math.~Z.}} \textbf{52} (1950),
 557--589.

\bibitem{T}
Tyc A., An elementary proof of the {W}eitzenb\"{o}ck theorem, \href{https://doi.org/10.4064/cm-78-1-123-132}{\textit{Colloq.
 Math.}} \textbf{78} (1998), 123--132.

\bibitem{U}
Umirbaev U.U., The {A}nick automorphism of free associative algebras,
 \href{https://doi.org/10.1515/CRELLE.2007.030}{\textit{J.~Reine Angew. Math.}} \textbf{605} (2007), 165--178,
 \href{https://arxiv.org/abs/math.RA/0607029}{arXiv:math.RA/0607029}.

\bibitem{E}
van~den Essen A., Polynomial automorphisms and the {J}acobian conjecture,
 \textit{Progress in Mathematics}, Vol.~190, \href{https://doi.org/10.1007/978-3-0348-8440-2}{Birkh\"{a}user Verlag}, Basel,
 2000.

\bibitem{K}
van~der Kulk W., On polynomial rings in two variables, \textit{Nieuw Arch.
 Wiskunde} \textbf{1} (1953), 33--41.

\bibitem{W}
Weitzenb\"{o}ck R., \"{U}ber die {I}nvarianten von linearen {G}ruppen,
 \href{https://doi.org/10.1007/BF02547779}{\textit{Acta Math.}} \textbf{58} (1932), 231--293.

\end{thebibliography}
